\section{Projective bundles over abelian torsors}
\label{sec:torsor-projective-bundles}

The following theorem excludes scheme projective bundles over
nonschematic abelian torsors over every base.

\begin{theorem}\label{thm:torsor-projective-bundle}
Let $S$ be an arbitrary scheme, let $A$ be an abelian $S$-scheme,
and let $Q$ be an fppf $A$-torsor algebraic space. Let $E$ be a finite
locally free $\OO_Q$-module whose rank is everywhere positive.
Then $\PP_Q(E)$ is a scheme if and only if $Q$ is a scheme.
\end{theorem}

An abelian scheme here is a smooth proper group scheme of finite
presentation with geometrically connected fibres; projectivity is not part of
the definition. We use the convention that $\PP_Q(E)$ parametrizes
invertible quotients, with tautological sheaf $\OO_{\PP_Q(E)}(1)$.
The bundle is locally free on the \'etale site of $Q$ and need not
have constant rank over the whole base. The converse implication
in the theorem is the usual construction of a projective bundle
over a scheme.

The proof begins with the normal local case. Affine boundaries and
the Ramanujam--Samuel theorem supply line bundles whose numerical
classes cannot vanish along a common positive-dimensional
translation kernel. For a general local base we apply this argument
simultaneously over its finite normalization. Selected relative
Picard classes provide projective targets, making that normalized
torsor projective. Raynaud's finite-cover criterion then descends
schematicity. Approximation removes the noetherian hypothesis.

Nonschematic elliptic torsors already occur over normal
noetherian local surfaces \cite[Theorem~2.1 and \S3]{Zomervrucht}.
The theorem applies to vector bundles of every positive rank on
these torsors, in every relative dimension, without a homogeneity
assumption. Brion's classification concerns homogeneous projective
bundles over abelian varieties over an algebraically closed field
\cite[pp.~2475--2479]{BrionProjective}.

\subsection{Affine boundaries and translation kernels}

\begin{lemma}\label{lem:normal-affine-relative-boundary}
Let $R$ be a normal noetherian local ring, with closed point
$s\in S=\Spec R$. Let $X\to S$ be a smooth proper scheme with
geometrically integral fibres. For every affine open $U\subset X$
meeting $X_s$, there is an effective Cartier divisor $D$ on $X$
whose support is $X\setminus U$. Its restriction to every fibre is
an effective Cartier divisor, and $D$ is flat over $S$.
\end{lemma}

\begin{proof}
The ring $R$ is a domain. Smoothness over the normal base makes $X$
normal. It is also irreducible. Indeed, a nonempty open of $X$ has
nonempty open image in the integral scheme $S$, so it meets the
generic fibre. Two such opens meet because that fibre is integral.
Thus $X$ is integral.

The map $U\to S$ is smooth and hence open. Its image contains the
closed point of the local scheme $S$, so its image is all of $S$.
Consequently $U_t$ is nonempty and dense in $X_t$ for every $t\in S$.
Since $X$ is separated over the affine scheme $S$, it is a separated
scheme. Its affine opens have affine intersections, and therefore
the inclusion $U\hookrightarrow X$ is an affine morphism.

By EGA IV$_4$, 21.12.7, every irreducible component of $X\setminus U$
has codimension one \cite[pp.~307--308]{EGA}. There are finitely many
such components. Let $D$ be the Weil cycle which is their sum, with
coefficient one. Its support contains no irreducible component of
any fibre, since $U_t$ is dense in the integral fibre $X_t$.
The second assertion of EGA IV$_4$, 21.14.3(i), applies to the smooth
morphism $X\to S$ and the normal locally noetherian base: precisely
such codimension-one cycles are locally principal
\cite[pp.~325--328]{EGA}.

Thus $D$ is a Cartier divisor. It is effective as a Cartier divisor:
a local rational equation has nonnegative valuation at every
codimension-one point, so normality makes it regular. Since $X$ is
integral, that equation is a non-zero-divisor. Its zero set is exactly
$X\setminus U$. On each integral fibre the equation remains nonzero,
because $U_t$ is nonempty; it is again a non-zero-divisor. Finally,
$\OO_X$ and $\OO_X(-D)$ are flat over $S$, and
\[
 0\longrightarrow\OO_X(-D)\longrightarrow\OO_X
 \longrightarrow\OO_D\longrightarrow0
\]
remains exact on every fibre. The fibre criterion for flatness gives
the flatness of $D\to S$.
\end{proof}

For a line bundle $L$ on an abelian variety $B$ over an algebraically
closed field, write
\[
 \phi_L:B\longrightarrow B^\vee,
 \qquad b\longmapsto[t_b^*L\otimes L^{-1}].
\]
The theorem of the square makes this a homomorphism. Its kernel is
all of $B$ exactly when $L$ lies in $\Pic^0(B)$, equivalently when
$L$ is algebraically trivial
\cite[Proposition~9.2, Remark~9.3 and Proposition~10.1,
pp.~14--17 of the author revision]{MilneAV}.
If $i:B'\hookrightarrow B$ is an abelian subvariety and
$\phi_L\circ i=0$, then
\[
 \phi_{i^*L}=i^\vee\circ\phi_L\circ i=0.
\]
It follows that $i^*L$ is algebraically, and hence numerically, trivial.
We will apply this implication to restrictions of global line
bundles on a torsor.

\begin{lemma}\label{lem:torsor-finite-translation-kernel}
Let $S$ be the spectrum of a noetherian local ring, with closed point
$s$, let $A/S$ be an abelian scheme, and let $Q$ be an fppf
$A$-torsor algebraic space. Choose a geometric point $\bar s$ above
$s$ and a point $q_0\in Q_{\bar s}$. For a line bundle $M$ on $Q$,
put $L_M=\tau_{q_0}^*(M|_{Q_{\bar s}})$, where
$\tau_{q_0}:A_{\bar s}\to Q_{\bar s}$ is the torsor isomorphism.
If there are finitely many line bundles $M_1,\ldots,M_n$ on $Q$ such
that
\[
 \bigcap_{i=1}^n\ker\phi_{L_{M_i}}\subset A_{\bar s}
\]
is finite, then $Q$ is a scheme.
\end{lemma}

\begin{proof}
Corollary~\ref{cor:picard-schematic-abelian-torsor}, applied to $A$
itself, shows that $P=\Pic_{A/S}$ is a separated $S$-scheme locally
of finite presentation. For every line bundle $M$ on $Q$, the action
map $A\times_S Q\to Q$ pulls $M$ back to a family of line bundles
on $A$, and therefore defines a morphism
\[
 f_M:Q\longrightarrow P.
\]
More explicitly, a point $q\in Q(T)$ is sent to the relative Picard
class of $\tau_q^*M$ on $A_T$. On the geometric closed fibre, the
identification given by $q_0$ writes $f_M$ as
\[
 a\longmapsto[L_M]+\phi_{L_M}(a).
\]
In particular the nonempty fibres of
$f=(f_{M_1},\ldots,f_{M_n}):Q\to P^n$ over $\bar s$ are translates
of the indicated finite intersection of kernels.

The morphism $f$ is proper. Its graph in $Q\times_S P^n$ is a closed
immersion because $P^n$ is separated over $S$, and the projection
$Q\times_S P^n\to P^n$ is proper because $Q\to S$ is proper.
It is quasi-finite at every point over $s$, as can be checked after
the residue-field extension to $\bar s$. The non-quasi-finite locus
is closed in $Q$. Its image in $S$ is closed by properness, and does
not contain $s$. A nonempty closed subset of a local scheme contains
its closed point, so this locus is empty. Hence $f$ is proper and
quasi-finite, and is finite \cite[Tag~0A4X]{Stacks}.
A finite algebraic space over a scheme is a scheme, proving the lemma.
\end{proof}

The lemma does not require a polarization, flatness of the finite
kernel, or local constancy of N\'eron--Severi groups. The use of the
full Picard scheme as target avoids choosing components in a family.

\subsection{The normal noetherian local case}

\begin{proposition}\label{prop:normal-local-torsor-projective-bundle}
Theorem~\ref{thm:torsor-projective-bundle} holds when $S$ is
the spectrum of a normal noetherian local ring.
\end{proposition}

\begin{proof}
Write $X=\PP_Q(E)$, let $\pi:X\to Q$ be the projection, and suppose
that $X$ is a scheme. The torsor $Q\to S$ is smooth, proper and
finitely presented, with geometrically integral fibres. It is
connected, since its fibres are connected and the base is local.
Thus $E$ has constant rank $r$. If $r=1$, the morphism $\pi$ is an
isomorphism and the conclusion follows. Assume $r\geq2$.
The scheme $X\to S$ is smooth and proper with geometrically integral
fibres, so Lemma~\ref{lem:normal-affine-relative-boundary} applies.

We first record the line-bundle decomposition on $X$. For every
invertible sheaf $H$ on $X$ there are an integer $m$ and an invertible
sheaf $M$ on $Q$ with
\begin{equation}\label{eq:torsor-bundle-picard-decomposition}
 H\simeq\OO_X(m)\otimes\pi^*M.
\end{equation}
This formula is valid for the algebraic-space base $Q$. On affine
\'etale charts of $Q$, the ordinary projective-bundle formula gives
a locally constant fibre degree. These degrees agree on overlaps,
and connectedness makes their value $m$ constant. After twisting
by $\OO_X(-m)$, every projective-space fibre has trivial restriction.
Cohomology and base change makes $\pi_*(H(-m))$ invertible, and the
evaluation map is an isomorphism on every fibre, hence an isomorphism.
The Hilbert-polynomial and base-change assertions used here are
EGA III$_2$, 7.9.11 and 7.9.8--7.9.9, together with 7.7.5
\cite{EGA}. These statements and the resulting line bundle
descend along the
\'etale charts, giving \eqref{eq:torsor-bundle-picard-decomposition}.
The same calculation proves uniqueness.

Let $U\subset X$ be an affine open meeting the closed fibre, and let
$D$ be its effective Cartier boundary from
Lemma~\ref{lem:normal-affine-relative-boundary}. In
\eqref{eq:torsor-bundle-picard-decomposition} for $H=\OO_X(D)$,
one has $m\geq0$. Indeed, choose a geometric point of $U$ and
restrict the canonical section of $\OO_X(D)$ to its
$\pi$-fibre $\PP^{r-1}$. This is a nonzero section of $\OO(m)$,
which is impossible for $m<0$. Other $\pi$-fibres may be contained
in $D$; this sign argument makes no flatness assertion over $Q$.

Suppose, towards a contradiction, that $Q$ is not a scheme. Put
$k=\overline{\kappa(s)}$, choose $q_0\in Q_k$, and use it to
identify $Q_k$ with $A_k$. For every global line bundle $M$ on $Q$
define $L_M$ as in Lemma~\ref{lem:torsor-finite-translation-kernel}.
Consider the scheme-theoretic intersection
\[
 K=\bigcap_{M\in\Pic(Q)}\ker\phi_{L_M}\subset A_k.
\]
This intersection is already a finite intersection. To see this,
its defining quasi-coherent ideal is the sum of the ideals of the
individual kernels. On the noetherian scheme $A_k$ this ideal is
of finite type. On each member of a finite affine cover, finitely
many sections drawn from finitely many of those ideals generate
the sum. Taking the union of the finite lists proves the assertion.
Thus no finite-generation assertion about $\Pic(Q)$ or its image
in a N\'eron--Severi group is needed.

By Lemma~\ref{lem:torsor-finite-translation-kernel}, $K$ is not
finite. A zero-dimensional closed subscheme of the proper
$k$-scheme $A_k$ is finite, so $K$ has positive dimension. Over the
algebraically closed field $k$, its reduced identity component
\[
 B=(K_{\mathrm{red}})^0
\]
is consequently a positive-dimensional abelian subvariety of $A_k$.
Choose an integral projective curve $C_0\subset B$ and let
$C=\tau_{q_0}(C_0)\subset Q_k$. Such a curve exists because an
abelian variety over a field is projective. Smoothness of the curve
is not needed.

For every global line bundle $M$ on $Q$, the homomorphism
$\phi_{L_M}$ vanishes on $B$. The observation preceding
Lemma~\ref{lem:torsor-finite-translation-kernel} shows that
$L_M|_B$ is algebraically trivial. In particular,
\begin{equation}\label{eq:torsor-curve-degrees-zero}
 \deg(M|_C)=0\quad\text{for every }M\in\Pic(Q),
 \qquad \deg((\det E)|_C)=0.
\end{equation}
Set $Y=\PP_C(E|_C)$, an integral projective $k$-scheme of dimension
$r$, and put $\xi=c_1(\OO_Y(1))$. The inclusion $C\hookrightarrow
Q_k$ identifies $Y$ with a closed subscheme of $X_k$.
The projective-bundle formula and Chern relation give
\begin{equation}\label{eq:torsor-bundle-top-intersection-zero}
 \deg\bigl(\xi^r\cap[Y]\bigr)
   =\deg((\det E)|_C)=0.
\end{equation}
For clarity, with the quotient convention the relation over the
curve is
$\xi^r\cap[Y]=\pi_C^*c_1(E|_C)\,\xi^{r-1}\cap[Y]$;
pushing forward uses
$\pi_{C*}(\xi^{r-1}\cap[Y])=[C]$.
These formulas apply to an integral, possibly singular curve
\cite[Tags~02TX and~02TZ, equation~(42.37.1.1)]{Stacks}.

We now prove, for every positive-dimensional integral closed
subvariety $Z\subset Y$, that
\begin{equation}\label{eq:torsor-bundle-positive-intersections}
 \deg\bigl(\xi^d\cap[Z]\bigr)>0,
 \qquad d=\dim Z.
\end{equation}
Choose a point of $Z$ and an affine open $U\subset X$ containing
its image. This open meets the closed fibre of $X\to S$.
Let $D$ be its boundary from
Lemma~\ref{lem:normal-affine-relative-boundary}, and write
\[
 \OO_X(D)\simeq\OO_X(m)\otimes\pi^*M,
 \qquad m\geq0.
\]
The mixed intersection contributed by $M$ is zero. More explicitly,
the projection formula gives
\[
 \deg\bigl(c_1(\pi_C^*(M|_C))\,\xi^{d-1}\cap[Z]\bigr)
 =\deg\bigl(c_1(M|_C)\cap
       \pi_{C*}(\xi^{d-1}\cap[Z])\bigr)=0.
\]
Indeed, the pushforward on the right is a one-cycle on the integral
curve $C$, hence an integer multiple of $[C]$, and
\eqref{eq:torsor-curve-degrees-zero} applies. This also verifies the
required numerical equality when $C$ is singular or $Z$ lies in
a fibre of $Y\to C$.

The open $U_k\cap Z$ is nonempty and affine: it is a closed
subscheme of the affine scheme $U_k$. It cannot be all of $Z$, since
a positive-dimensional integral projective scheme over a field is
not affine. Thus the restriction $D_Z$ is a nonempty effective
Cartier divisor on $Z$. Its associated cycle has the form
\[
 [D_Z]=\sum_j a_j[Z_j],\qquad a_j>0,
\]
where every $Z_j$ is integral of dimension $d-1$.
If $d=1$, its degree is strictly positive, and the preceding
projection calculation yields
\[
 m\deg\bigl(\xi\cap[Z]\bigr)=\deg[D_Z]>0.
\]
Because $m\geq0$, this proves $m>0$ and
\eqref{eq:torsor-bundle-positive-intersections} in dimension one.

For $d>1$, assume the assertion in all smaller positive dimensions.
Intersecting $[D_Z]$ with $\xi^{d-1}$ gives
\[
 \begin{split}
 m\deg\bigl(\xi^d\cap[Z]\bigr)
  &=\deg\bigl(\xi^{d-1}\cap[D_Z]\bigr)\\
  &=\sum_j a_j\deg\bigl(\xi^{d-1}\cap[Z_j]\bigr)>0.
 \end{split}
\]
Again $m\geq0$ proves the desired strict positivity. The induction
is complete. Taking $Z=Y$ contradicts
\eqref{eq:torsor-bundle-top-intersection-zero}.
Therefore $Q$ is a scheme.
\end{proof}

The argument uses a boundary chosen separately for each integral
subvariety $Z$. A single effective boundary need not be ample.
What matters is that every such boundary has, on $Y$, its numerical
class on the same nonnegative ray, and that an appropriate boundary
cuts every positive-dimensional $Z$ nontrivially. The calculation
after extension to $k$ is entirely an intersection calculation;
schematicity is not descended from an algebraic closure of the
residue field.

\subsection{Projective targets for relative Picard classes}

We will need projectivity over a finite normal cover of a local
base. The following argument verifies the simultaneous affine
neighborhood needed for the group scheme.

\begin{lemma}\label{lem:dual-projective-finite-normal-cover}
Let $S$ be noetherian local, let $W\to S$ be finite and surjective
with $W$ normal, and let $B$ be an abelian $S$-scheme. Then
$B_W$ is projective over $W$.
\end{lemma}

\begin{proof}
Choose an affine open $V\subset B$ meeting the closed fibre.
Since $V\to S$ is smooth and open, its image is all of $S$.
Translation therefore saturates $V$ to $B$, as can be checked on
geometric fibres. The hypotheses of \cite[6.6, Theorem~2(c)]{BLR}
are satisfied: the base is affine and locally noetherian, the group
is smooth with connected fibres, and the saturated open is
quasi-affine. Hence every finite set in $B$ has a common affine
open neighborhood.

The map $B_W\to B$ is finite. Thus the same property holds in $B_W$:
put the finitely many image points in an affine open of $B$ and
take its affine inverse image. The scheme $W$ is affine,
noetherian and semilocal. Choose one point in every closed
$W$-fibre of $B_W$, and put these finitely many points in a single
affine open $V_W\subset B_W$. Its image in $W$ is open by
smoothness and contains all closed points, so is all of $W$.
Indeed a nonempty closed subset of an affine scheme contains a
closed point. Translation saturates $V_W$ to $B_W$.
Since $W$ is normal, \cite[6.6, Theorem~2(d)]{BLR} gives a
$W$-ample invertible sheaf on $B_W$. Properness makes $B_W$
projective.
\end{proof}

\begin{lemma}\label{lem:projective-picard-class-torsor}
Let $T$ be a noetherian affine scheme and let $C/T$ be an abelian
scheme whose dual $B=C^\vee$ is projective over $T$. Put
$P=\Pic_{C/T}$, and let $\delta$ be a section of the fppf quotient
$P/B$ over $T$. Then
\[
 P_\delta=P\times_{P/B,\delta}T
\]
is a projective $T$-scheme.
\end{lemma}

\begin{proof}
The sheaf $P_\delta$ is a $B$-torsor. In particular it is an
algebraic space, smooth and proper of finite presentation over $T$.
No representability or local constancy assertion about the whole
quotient $P/B$ is needed. The Picard identities
\cite[\S\S3.1--3.2]{PR} say that
$\iota=[-1]^*$ acts as $-1$ on $B$ and as $+1$ on $P/B$.
Thus it preserves $P_\delta$. If $x$ is a local section of this
torsor and $b$ is a section of $B$, then
\[
 \iota(x+b)-(x+b)=\iota(x)-x-2b.
\]
Multiplication by $2$ on $B$ is finite locally free and surjective,
including in characteristic $2$
\cite[Corollary~4.6]{LaurentSchroer}.
The fixed subsheaf $R_\delta$ is consequently a torsor under
$B[2]$. By descent of finite locally free affine morphisms,
$R_\delta\to T$ is a finite locally free surjective scheme. On
a relative-dimension-$g$ component its degree is $2^{2g}$.

The inclusion $R_\delta\to P_\delta$ supplies a tautological
section after base change, and hence an isomorphism
\[
 (P_\delta)_{R_\delta}\simeq B_{R_\delta}.
\]
Only a relative Picard class is used here; a universal line bundle
on $C\times_TP_\delta$ is not required.
The right side is projective over the affine scheme $R_\delta$,
and therefore every finite set in it has an affine neighborhood.
Apply \cite[Theorem~4.2]{RaynaudDescent} with both the fpqc cover
and the separate finite surjective cover equal to
$R_\delta\to T$. The pullback is a locally noetherian scheme
with the required affine neighborhoods. It follows that
$P_\delta$ is a scheme.

We can now apply the norm to schemes. Write
$g:(P_\delta)_{R_\delta}\to P_\delta$ for the finite locally free
projection, and pull a $T$-ample invertible sheaf on $B$ back to
$B_{R_\delta}\simeq(P_\delta)_{R_\delta}$. Denote the resulting
sheaf by $L'$. It is ample over $T$, since $B_{R_\delta}\to B$ is
finite. Its norm is the invertible sheaf
\[
 N_g(L')=
 \det(g_*L')\otimes
 \det\bigl(g_*\OO_{(P_\delta)_{R_\delta}}\bigr)^{-1}.
\]
EGA II, 6.6.1, makes this sheaf ample over $T$
\cite{EGA}. Since $P_\delta$ is proper over $T$, it is projective.
\end{proof}

The torsor $P_\delta$ in this lemma need not be trivial.
Its finite flat symmetric-representative cover is essential to the
argument, and need not be \'etale. Likewise, an arbitrary schematic
abelian torsor over a normal semilocal base need not be projective:
\cite[Example~6.1, p.~1574]{BhattAnnihilation} gives a Zariski
locally trivial torsor of infinite order which cannot be split by
a finite cover. A projective smooth torsor over that semilocal base
would have a finite locally free multisection by
Lemma~\ref{lem:semilocal-multisection}. We next obtain the
additional projectivity needed here from the original scheme
projective bundle.

\subsection{Uniform boundaries over a finite normalization}

\begin{proposition}\label{prop:torsor-projective-after-normalization}
Let $S$ be the spectrum of an excellent noetherian local ring,
let $A/S$ be an abelian scheme, and let $Q$ be an fppf $A$-torsor
algebraic space. Suppose that $E$ is finite locally free of
everywhere positive rank and $X=\PP_Q(E)$ is a scheme. If
$W\to S$ is the normalization of $S_{\mathrm{red}}$, then
$Q_W$ is projective over $W$.
\end{proposition}

\begin{proof}
Let $s$ be the closed point of $S$. The morphism $W\to S$ is
finite and surjective by excellence. The scheme $W$ is normal,
noetherian, affine and semilocal, and all its closed points lie
over $s$. Its finitely many integral components are open and
closed. We do not assume that $W\to S$ is flat.
The dual $B=A^\vee$ is an abelian $S$-scheme and commutes with
base change \cite[\S3.1]{PR}\cite[Corollary~4.6]{LaurentSchroer}.
Lemma~\ref{lem:dual-projective-finite-normal-cover} makes $B_W$
projective over $W$.

The rank $r$ of $E$ is constant. Its locally constant rank
function is constant on each connected geometric fibre of
$Q\to S$, hence descends to the connected local base.
If $r=1$, then $X=Q$ is a scheme. An affine open in $Q$ meeting
the closed fibre is saturated by the $A$-action, so
\cite[6.6, Theorem~2(c)]{BLR} makes $Q$ AF. This property passes
through the finite map $Q_W\to Q$. One affine open in $Q_W$
therefore meets all closed $W$-fibres, and part (d) of the same
theorem makes $Q_W$ projective, exactly as in
Lemma~\ref{lem:dual-projective-finite-normal-cover}.
Assume henceforth that $r\geq2$.

\emph{Global boundaries.}
Let $U\subset X$ be an affine open meeting $X_s$. The map
$U\to S$ is smooth and open, so is surjective. Its finite
pullback $U_W$ is affine and meets every fibre over $W$.
Those fibres are geometrically integral, so $U_W$ is dense
in every fibre, including after a residue-field extension.
The scheme $X_W$ is smooth over normal $W$, hence normal;
on every integral component of $W$ it is integral, by the
generic-fibre argument in
Lemma~\ref{lem:normal-affine-relative-boundary}.

The inclusion $U_W\hookrightarrow X_W$ is affine because
$X_W$ is separated over affine $W$. EGA IV$_4$, 21.12.7,
therefore makes its boundary a finite union of prime divisors.
Give each component multiplicity one. The support of this
effective cycle contains no irreducible component of any
fibre, so EGA IV$_4$, 21.14.3(i), makes it locally principal
\cite{EGA}. Normality makes the local equations regular
non-zero-divisors, and equations with the same Weil divisor
differ by a unit. We obtain a global effective Cartier divisor
$D_U$ on $X_W$ with support $X_W\setminus U_W$.
It remains effective Cartier on every geometric fibre,
because the section is nonzero on the nonempty open
$(U_W)_{\bar w}$ of the integral fibre.

On each open-and-closed component of $W$, the chartwise
argument proving \eqref{eq:torsor-bundle-picard-decomposition} gives
\[
 \OO_{X_W}(D_U)\simeq
 \OO_{X_W}(m_U)\otimes\pi_W^*M_U.
\]
Here $M_U$ is a global line bundle on $Q_W$, and $m_U$ is a
locally constant integer function on $W$. The componentwise
line bundles glue over the disjoint union. Every value of
$m_U$ is nonnegative: on a projective-space fibre containing
a point of $U_W$, the canonical section restricts to a nonzero
section of $\OO(m_U)$. These boundaries and their factors
$M_U$ are defined over the whole of $W$.

\emph{Finite common kernels on every closed fibre.}
For every $M\in\Pic(Q_W)$ let
$\phi_M:A_W\to B_W$ be its translation homomorphism.
It is defined after a local choice of torsor point and is
independent of that choice \cite[\S3.6]{PR}.
For a closed $w\in W$, choose
$k=\overline{\kappa(w)}$ and put
\[
 K_w=\bigcap_{M\in\Pic(Q_W)}\ker(\phi_{M,k})\subset A_k.
\]
This is a finite intersection, by the noetherian ideal argument
in the proof of
Proposition~\ref{prop:normal-local-torsor-projective-bundle}.
We claim it is finite.

If not, its reduced identity component is a positive-dimensional
abelian subvariety. Choose a translated integral curve
$C\subset Q_k$ in that component, as in that proposition.
Every global line bundle on $Q_W$ has degree zero on $C$,
including $\det E_W$. Set
\[
 Y=\PP_C(E_W|_C),\qquad \xi=c_1(\OO_Y(1)).
\]
Then $Y$ is integral projective of dimension $r$, is closed
in $X_k$, and the Chern relation gives
\[
 \deg(\xi^r\cap[Y])=0.
\]

For completeness, the boundaries just constructed give the
contradictory positivity uniformly over $W$.
For a positive-dimensional integral $Z\subset Y$, choose
an affine open $U$ in the original scheme $X$ containing the
image of a point of $Z$. That image lies over $s$, so $U$
belongs to the preceding construction. The intersection
$(U_W)_k\cap Z$ is nonempty and affine, but cannot equal the
positive-dimensional projective $Z$.
Thus $D_U$ restricts to a nonempty effective Cartier divisor
$D_Z$ on $Z$. If $d=\dim Z$ and $m=m_U(w)\geq0$, projection
formula gives
\begin{equation}\label{eq:uniform-boundary-intersection}
 m\deg(\xi^d\cap[Z])=
 \deg(\xi^{d-1}\cap[D_Z]).
\end{equation}
Indeed the omitted term is
\[
 \deg\bigl(c_1(M_U|_C)\cap
       \pi_{C*}(\xi^{d-1}\cap[Z])\bigr)=0:
\]
the pushforward is a one-cycle on the integral curve $C$,
and $M_U|_C$ has degree zero.

For $d=1$, the right side of
\eqref{eq:uniform-boundary-intersection} is strictly positive.
For $d>1$, write
$[D_Z]=\sum_j a_j[Z_j]$ with $a_j>0$ and $\dim Z_j=d-1$.
Induction makes every degree
$\deg(\xi^{d-1}\cap[Z_j])$ positive. In either case
\eqref{eq:uniform-boundary-intersection}, together with $m\geq0$,
gives $\deg(\xi^d\cap[Z])>0$.
Taking $Z=Y$ contradicts its zero top intersection.
This proves that every $K_w$ is finite.

\emph{A finite map to projective Picard-class torsors.}
There are finitely many closed points in $W$. At each one,
select finitely many global line bundles defining the
finite intersection $K_w$, and take the union of these finite
selections. Denote the resulting bundles by $M_1,\ldots,M_n$.
Their common kernel is finite on every closed geometric fibre.

For each $M_i$, the action defines
\[
 f_i:Q_W\longrightarrow\Pic_{A_W/W},
 \qquad q\longmapsto[\tau_q^*M_i].
\]
Its class modulo $B_W$ is a single section $\delta_i$ over $W$.
Indeed local choices of $q$ differ by translation, and the
corresponding Picard classes differ by an element of $B_W$.
They therefore agree in the quotient and descend.
The associated homomorphism is $\phi_{M_i}$.
Thus $f_i$ factors through the torsor $P_{\delta_i}$.
Lemma~\ref{lem:projective-picard-class-torsor} makes each of
these torsors projective over $W$.

Consider
\[
 f=(f_1,\ldots,f_n):Q_W\longrightarrow
       \prod_{i=1}^n P_{\delta_i}.
\]
It is proper because its target is separated over $W$ and
$Q_W$ is proper over $W$. On each closed geometric fibre it is,
after a torsor-point choice, a translate of the product
homomorphism $(\phi_{M_i})_i$. Its nonempty fibres are
translates of the finite common kernel.
The non-quasi-finite locus of $f$ is closed in $Q_W$, and its
proper image in $W$ contains no closed point. It is therefore
empty. Hence $f$ is proper and quasi-finite, and so finite
\cite[Tag~0A4X]{Stacks}.
A finite algebraic space over its scheme target is a scheme.
The finite pullback of a $W$-ample invertible sheaf on that
target is $W$-ample on $Q_W$. Thus $Q_W$ is projective.

If $A$ has relative dimension zero, the common kernel is
already finite without choosing line bundles; take $n=0$
and the target to be $W$. The same finiteness argument applies.
\end{proof}

The affine opens in this proof come from $X$ over the original
local base. Their finite pullbacks meet every closed fibre of
$W$ at once. This is what makes their line bundles global over
the semilocal normalization and permits one finite family of
translation homomorphisms.

\subsection{Descent and approximation over arbitrary bases}

\begin{proof}[Proof of Theorem~\ref{thm:torsor-projective-bundle}]
First let $S$ be excellent noetherian local, with closed point
$s$, and suppose $X=\PP_Q(E)$ is a scheme. Let $W\to S$ be the
finite normalization of $S_{\mathrm{red}}$.
Proposition~\ref{prop:torsor-projective-after-normalization}
makes $Q_W$ projective over the affine $W$, so every finite
set of its points has a common affine open neighborhood.

Choose an affine open $U\subset X$ meeting $X_s$.
The smooth map $U\to S$ is surjective, hence is an fpqc cover.
The morphism $U\to X\to Q$ gives a section of $Q_U$, so
$Q_U\simeq A_U$ is a locally noetherian scheme.
Apply \cite[Theorem~4.2]{RaynaudDescent} to the fppf sheaf $Q$,
using $U\to S$ as the fpqc cover and $W\to S$ as the separate
finite surjective cover. All finite orbits associated to
$W\times_S W$ lie in affine opens of $Q_W$. It follows that
$Q$ is a scheme. The normalization was not assumed flat;
the original nilpotent structure of $S$ is present in
$U\to S$.

Now let $S$ be excellent and noetherian. For $s\in S$, the
local result makes $Q\times_S\Spec\OO_{S,s}$ a scheme.
We spread a finite affine open cover of it to a neighborhood
of $s$. Choose an affine neighborhood $S_0$.
The finitely many affine charts and their quasi-compact open
immersions into the localized $Q$ are of finite presentation.
EGA IV$_3$, 8.8.2 and 8.10.5(iii), together with
\cite[Tag~07SK]{Stacks}, descend them to open immersions over
an affine neighborhood $T\subset S_0$ of $s$.
For an algebraic-space target the open-immersion assertion is
checked on a finite affine \'etale presentation and descends
from those charts. The complement of their finite union is
closed in $Q_T$ and has empty fibre over $s$. Its image in
$T$ is closed by properness. Removing that image gives a
neighborhood on which the charts cover, so $Q$ is a scheme
there. These neighborhoods cover $S$.

Finally take an arbitrary affine base $S=\Spec R$ and write
$R=\varinjlim R_i$ as the filtered union of its finitely
generated $\Z$-subalgebras. Each $R_i$ is excellent and
noetherian. We descend all the finite-presentation data needed
for a model of the hypothesis.
Descend the schemes $A$ and $X$, the algebraic space $Q$,
the group and action maps, and the isomorphism
\[
 A\times_R Q\longrightarrow Q\times_R Q,\qquad
 (a,q)\longmapsto(a+q,q),
\]
with its inverse and both composition identities.
Choose a finite affine \'etale presentation $V\to Q$
on which $E$ is trivial on each chart. Descend that
presentation, its overlaps, the vector-bundle transition
matrices and their inverse and cocycle identities.
This produces a finite locally free $E_i$ of positive ranks.
Also descend
\[
 X\simeq\PP_Q(E),
\]
its inverse and its two identity equations.
Thus $X_i$ is a scheme and is isomorphic to
$\PP_{Q_i}(E_i)$.
These are the scheme and algebraic-space limit theorems
EGA IV$_3$, \S\S8.8--8.10, and \cite[Tag~07SK]{Stacks}.

After increasing $i$, properness of $A_i$ follows from
EGA IV$_3$, 8.10.5(xii), and smoothness by EGA IV$_4$,
17.7.8(ii). The group identities hold at the finite stage.
The geometric-integrality failure locus of its fibres is
constructible by EGA IV$_3$, 9.7.7.
Its inverse image in $\Spec R$ is empty and the property
is preserved and reflected by residue-field extension.
Hence that constructible locus disappears at some later
stage. To see this explicitly, decompose it into finitely
many locally closed affine pieces; emptiness after the
filtered colimit means that $1$ belongs to their defining
ideals after localization, a finite relation already true
at a finite stage. Thus $A_i$ is an abelian scheme.

We also arrange \'etaleness and surjectivity of the presentation, and
smoothness and surjectivity of $V_i\to\Spec R_i$, using
EGA IV$_4$, 17.7.8(ii), and IV$_3$, 8.10.5(vi).
The map $V_i\to Q_i$ supplies local sections. Together with
the descended simple-transitivity isomorphism, it makes
$Q_i$ an fppf $A_i$-torsor.
Its smoothness and properness follow by fppf descent from
$A_i$. The construction therefore gives an fppf torsor,
a vector bundle of positive rank, and a scheme projective
bundle over the excellent noetherian model.
The case already proved makes $Q_i$ a scheme, and base
change makes $Q$ a scheme. Covering an arbitrary $S$ by
affines proves the theorem.
\end{proof}

The theorem excludes the projective-bundle construction over
nonschematic abelian torsors for all bases. Schematicity of the
Picard functor of a general proper flat family with geometrically
normal or geometrically integral fibres remains a separate
question: such a family need not have the torsor action and
translation homomorphisms used above.
